\chapter{系统测试与结果分析}

为了验证本系统的功能和性能是否满足设计要求，本章从测试环境搭建、功能测试、性能测试、分类精度测试以及用户体验测试五个方面对系统进行全面测试，并对测试结果进行分析总结。

\section{测试环境}

本系统的测试环境分为硬件环境和软件环境两部分。硬件环境主要采用华为Mate 60 Pro手机作为测试设备，其详细配置如表\ref{tab:hw-env}所示。

\begin{table}[htbp]
\centering
\caption{硬件测试环境}
\label{tab:hw-env}
\begin{tabular}{ll}
\toprule
项目 & 参数 \\
\midrule
设备型号 & 华为 Mate 60 Pro \\
处理器 & 麒麟 9000S \\
运行内存 & 12\,GB \\
存储空间 & 256\,GB \\
操作系统 & HarmonyOS 4.0 \\
屏幕尺寸 & 6.82 英寸 \\
电池容量 & 5000\,mAh \\
\bottomrule
\end{tabular}
\end{table}

软件测试环境如表\ref{tab:sw-env}所示。

\begin{table}[htbp]
\centering
\caption{软件测试环境}
\label{tab:sw-env}
\begin{tabular}{ll}
\toprule
项目 & 版本/说明 \\
\midrule
开发环境 & DevEco Studio 4.1 \\
SDK版本 & HarmonyOS SDK API 10 \\
推理框架 & MindSpore Lite 2.2 \\
模型格式 & .ms（MindSpore Lite） \\
编程语言 & ArkTS / eTS \\
测试框架 & HvTestUnit \\
\bottomrule
\end{tabular}
\end{table}

\section{功能测试}

功能测试旨在验证系统各功能模块是否按照需求规格说明正确运行。本文针对图片选择、图像分类推理、结果展示、历史记录保存和离线模式五个核心功能设计了测试用例，测试结果如表\ref{tab:func-test}所示。

\begin{table}[htbp]
\centering
\caption{功能测试用例及结果}
\label{tab:func-test}
\begin{tabular}{cllllc}
\toprule
编号 & 测试内容 & 预期结果 & 实际结果 & 结果 \\
\midrule
F-01 & 从相册选择一张图片 & 成功加载图片并显示预览 & 成功加载图片并显示预览 & 通过 \\
F-02 & 使用相机拍摄一张图片 & 成功拍照并加载图片预览 & 成功拍照并加载图片预览 & 通过 \\
F-03 & 对选中图片执行分类推理 & 返回Top-5分类结果及置信度 & 返回Top-5分类结果及置信度 & 通过 \\
F-04 & 分类结果界面展示 & 显示类别标签、置信度柱状图 & 显示类别标签、置信度柱状图 & 通过 \\
F-05 & 分类结果保存至历史记录 & 结果自动保存至本地历史列表 & 结果自动保存至本地历史列表 & 通过 \\
F-06 & 查看历史分类记录 & 列表展示所有历史记录及缩略图 & 列表展示所有历史记录及缩略图 & 通过 \\
F-07 & 断网情况下执行分类推理 & 离线模型正常推理并返回结果 & 离线模型正常推理并返回结果 & 通过 \\
F-08 & 连续进行10次分类推理 & 每次均正确返回分类结果 & 每次均正确返回分类结果 & 通过 \\
\bottomrule
\end{tabular}
\end{table}

由表\ref{tab:func-test}可知，所有8项功能测试用例均通过，表明系统的核心功能模块工作正常，满足设计需求。其中F-07测试在关闭Wi-Fi和移动数据的条件下进行，验证了系统基于MindSpore Lite端侧推理的离线可用性。

\section{性能测试}

性能测试主要从推理速度、内存占用和电池消耗三个维度评估系统在实际设备上的运行效率。

\subsection{推理时间测试}

选取100张不同类别和分辨率的图片，依次在测试设备上进行分类推理，记录每次推理耗时。推理时间统计结果如表\ref{tab:infer-time}所示。

\begin{table}[htbp]
\centering
\caption{推理时间测试结果}
\label{tab:infer-time}
\begin{tabular}{ll}
\toprule
统计指标 & 推理时间(ms) \\
\midrule
平均推理时间 & 45 \\
最小推理时间 & 32 \\
最大推理时间 & 78 \\
中位数 & 42 \\
标准差 & 8.3 \\
\bottomrule
\end{tabular}
\end{table}

从表\ref{tab:infer-time}可以看出，系统的平均推理时间为45\,ms，满足实时分类的需求。推理时间的波动主要受图片分辨率差异影响——高分辨率图片在预处理阶段的缩放操作会增加少量耗时。

\subsection{内存占用测试}

在系统运行过程中，通过DevEco Studio Profiler监控应用的内存使用情况，测试结果如表\ref{tab:mem-test}所示。

\begin{table}[htbp]
\centering
\caption{内存占用测试结果}
\label{tab:mem-test}
\begin{tabular}{ll}
\toprule
统计指标 & 内存占用(MB) \\
\midrule
平均内存占用 & 120 \\
峰值内存占用 & 156 \\
空闲状态内存 & 68 \\
\bottomrule
\end{tabular}
\end{table}

系统在执行推理时，内存从空闲状态的68\,MB上升至平均120\,MB，峰值156\,MB。峰值出现在模型加载与图片预处理并行的时刻，之后内存回落至稳定水平，未出现内存泄漏现象。

\subsection{电池消耗测试}

在设备充满电的情况下，连续执行100次分类推理操作（含图片选择和结果展示），测试前后电池电量的变化。测试结果显示，100次连续推理消耗电池电量约2\%，折算单次推理耗电约0.02\%，表明系统的能耗表现良好，适合移动端长时间使用。

\subsection{模型版本对比测试}

为评估模型优化策略的效果，对原始FP32模型、量化INT8模型和蒸馏MobileNetV2模型三个版本进行对比测试，结果如表\ref{tab:model-compare}所示。

\begin{table}[htbp]
\centering
\caption{不同模型版本对比结果}
\label{tab:model-compare}
\begin{tabular}{lccc}
\toprule
模型版本 & 准确率(\%) & 模型大小(MB) & 推理时间(ms) \\
\midrule
原始FP32模型 & 93.5 & 45.2 & 112 \\
量化INT8模型 & 92.3 & 11.8 & 45 \\
蒸馏MobileNetV2模型 & 89.7 & 8.6 & 38 \\
\bottomrule
\end{tabular}
\end{table}

从表\ref{tab:model-compare}可以看出，量化INT8模型在准确率仅下降1.2个百分点的情况下，模型大小缩减为原始模型的26.1\%，推理速度提升了约2.5倍，是准确率与效率之间的最佳平衡点，因此本系统最终选用量化INT8模型作为端侧部署版本。

\section{分类精度测试}

\subsection{总体精度}

为全面评估模型的分类性能，在CIFAR-10测试集（共10000张图片）上对量化INT8模型进行分类精度测试。测试结果表明，模型在CIFAR-10测试集上的总体准确率为92.3\%。

\subsection{各类别精度}

模型在CIFAR-10各类别上的分类精度如表\ref{tab:cls-acc}所示。

\begin{table}[htbp]
\centering
\caption{CIFAR-10各类别分类精度}
\label{tab:cls-acc}
\begin{tabular}{lcc}
\toprule
类别 & 英文名称 & 准确率(\%) \\
\midrule
飞机 & Airplane & 94.2 \\
汽车 & Automobile & 96.1 \\
鸟 & Bird & 88.7 \\
猫 & Cat & 85.3 \\
鹿 & Deer & 91.8 \\
狗 & Dog & 87.5 \\
青蛙 & Frog & 93.6 \\
马 & Horse & 94.8 \\
船 & Ship & 95.2 \\
卡车 & Truck & 95.8 \\
\bottomrule
\end{tabular}
\end{table}

需要说明的是，表\ref{tab:cls-acc}中各类别准确率的计算方式为该类别正确分类的样本数占该类别总样本数的比例。从测试结果来看，汽车类别的分类精度最高，达到96.1\%；猫类别的分类精度最低，为85.3\%。猫和狗两个类别的精度相对较低，这与猫和狗在视觉特征上存在较高相似性有关——在混淆矩阵中可以观察到猫和狗之间存在一定程度的误分类现象。

\subsection{混淆矩阵分析}

为进一步分析分类错误的具体分布情况，本文对CIFAR-10测试集的预测结果构建了混淆矩阵。混淆矩阵的分析结果显示以下主要特征：

（1）猫(Cat)和狗(Dog)之间存在较明显的混淆，约7.2\%的猫被误分类为狗，5.8\%的狗被误分类为猫，这与两类动物在形态、毛色等视觉特征上的相似性有关。

（2）鸟(Bird)和鹿(Deer)之间偶有混淆，主要由于部分自然环境背景的干扰导致特征提取出现偏差。

（3）飞机(Airplane)和船(Ship)之间存在少量混淆，两者在天空和水面的简单背景下存在一定的形状相似性。

总体而言，模型在大多数类别上表现良好，混淆主要集中在视觉特征相似的类别之间，符合预期。

\section{用户体验测试}

\subsection{测试方法}

邀请20名不同年龄和职业背景的用户参与体验测试，测试流程如下：首先向用户简要介绍系统功能，然后让用户自主完成图片选择、分类推理、结果查看和记录回顾等操作，最后请用户填写系统可用性量表（System Usability Scale, SUS）和满意度调查问卷。

\subsection{SUS评分}

SUS量表包含10个题目，采用5级评分制，最终得分范围为0--100分。根据20名用户的评分数据，系统SUS得分为82.5分。根据SUS评分等级划分，82.5分属于"优秀"等级（A级），表明系统的可用性总体表现良好。

\subsection{用户满意度调查}

用户满意度调查从五个维度进行评价，各维度的平均得分如表\ref{tab:satisfaction}所示。

\begin{table}[htbp]
\centering
\caption{用户满意度调查结果}
\label{tab:satisfaction}
\begin{tabular}{lcc}
\toprule
评价维度 & 满分 & 平均得分 \\
\midrule
界面美观度 & 5 & 4.3 \\
操作便捷性 & 5 & 4.5 \\
分类准确性 & 5 & 4.1 \\
响应速度 & 5 & 4.6 \\
总体满意度 & 5 & 4.3 \\
\bottomrule
\end{tabular}
\end{table}

从表\ref{tab:satisfaction}可以看出，用户对"响应速度"的评价最高（4.6分），这与端侧推理带来的低延迟体验直接相关；对"分类准确性"的评价相对较低（4.1分），主要原因是猫、狗等相似类别的误分类被用户感知到，后续可通过增加训练数据或采用更精细的模型结构来改善。

\section{测试结果分析}

本章对系统进行了功能、性能、分类精度和用户体验四个方面的测试，现对测试结果进行综合分析。

\textbf{（1）功能完整性。}8项功能测试用例全部通过，表明系统在图片选择、分类推理、结果展示、历史保存和离线模式等方面的功能实现完整且稳定，满足需求规格说明的要求。

\textbf{（2）性能优化效果。}通过INT8量化技术，模型大小从45.2\,MB压缩至11.8\,MB，压缩比达73.9\%；推理时间从112\,ms缩短至45\,ms，加速比为2.5倍，实现了显著的模型轻量化与推理加速，使得端侧实时分类成为可能。同时，量化带来的准确率损失仅为1.2个百分点（从93.5\%降至92.3\%），在可接受的范围内，验证了量化策略的有效性。

\textbf{（3）分类精度分析。}模型在CIFAR-10测试集上的总体准确率为92.3\%，各类别精度介于85.3\%至96.1\%之间。猫和狗两个类别精度偏低，混淆矩阵分析表明二者存在交叉误分类现象，其根本原因在于细粒度视觉特征的区分难度较大。未来可通过引入注意力机制或增加细粒度类别的训练数据来提升这些类别的识别精度。

\textbf{（4）用户体验。}系统SUS得分为82.5分，用户满意度调查平均得分4.3分（满分5分），整体处于"优秀"水平。响应速度维度得分最高（4.6分），充分体现了鸿蒙平台与MindSpore Lite端侧推理的低延迟优势；分类准确性维度得分相对偏低（4.1分），与猫狗等细粒度类别的误分类直接相关，是后续优化需重点关注的方面。

综上所述，本系统在功能完整性、运行性能、分类精度和用户体验方面均达到了预期目标。量化INT8模型在准确率与推理效率之间取得了较好的平衡，适合鸿蒙移动终端的实际部署需求。后续工作将重点关注细粒度类别的精度提升和更多应用场景的适配。